% =============================================================================
%  «ТЕОРИЯ СИСТЕМ» — ствол серии «Архитектура Рассуждения»
%  Блок 7. ГРАНИЦА: одна несущая стена (КУЛЬМИНАЦИЯ).  Разделы 7.1–7.5 (черновик прозы).
%  Стиль: «выводим и показываем»; репо-якоря — в приложении, не в сносках.
%  Самокомпилируемо: pdflatex blok-7-granica.tex (см. память latex-build-windows).
% =============================================================================
\documentclass[12pt,a4paper]{article}
\usepackage{cmap}
\usepackage[T2A]{fontenc}
\usepackage[utf8]{inputenc}
\usepackage[russian]{babel}
\usepackage{paratype}
\usepackage{amsmath,amssymb,amsthm}
\usepackage{listings}
\usepackage[expansion=false]{microtype}
\usepackage[margin=28mm]{geometry}

\newcommand{\rezultat}{\medskip\noindent\emph{Что отсюда следует.}\ }

\title{\textbf{Теория Систем}\\[2mm]\large Блок 7. Граница: одна несущая стена}
\author{}
\date{}

\begin{document}
\maketitle
\thispagestyle{empty}

\noindent\small\textit{Черновик: разделы 7.1--7.5 (кульминация). Продолжает Блок~6 (вертикаль).
Объект, время, целое и вертикаль построены --- здесь находится их \emph{единственная} граница, и
она внутри собственного ядра теории. Машинные якоря --- в Приложении.}\normalsize

\bigskip

\noindent Теория построена. Спросим теперь, где её \emph{единственная} несущая стена. Ответ собирает
обе сквозные ноты: позитивное начало (единица, M1) и предел (role-limit) --- это \emph{две стороны
одной} L4-асимметрии. И главное: граница, которую теория предсказывает в мире, обнаружится
\emph{внутри неё самой}.

% =============================================================================
\section*{7.1\quad Топос-вопрос: характеризация}
% =============================================================================

{\sloppy
У зрелой категории спрашивают: есть ли у неё классификатор подобъектов --- то, что делает её
\emph{топосом}? Ответ здесь --- не «да» и не «нет», а \emph{характеризация}: подобъект классифицируется
\emph{ровно тогда}, когда задающий его предикат \emph{разрешим} и \emph{Roles-насыщен}.

Две стены складываются в одну карту. Классифицируемость влечёт разрешимость --- а это в точности
H1; и Roles-насыщенность --- а это ярус Roles. Стало быть, провал полной топосности --- не глухая
стена, а именно H1 плюс ограничение яруса Roles.\par}

\rezultat
Топос-вопрос решён \emph{картой}, а не стеной: видно, \emph{где} категория ведёт себя как топос и
\emph{почему} за этой границей --- нет. Честно: это \textbf{не} «теория систем есть топос»; топос ---
лишь на разрешимом Roles-насыщенном фрагменте.

% =============================================================================
\section*{7.2\quad Мост к процессам}
% =============================================================================

{\sloppy
Статичный триад и динамический процесс --- \emph{две грани одной субстанции}. Динамика системы
(Блок~4) \emph{есть} процесс; рациональная траектория \emph{есть} вещественное-как-процесс
(например, $(1/2)^n$ --- геометрическое половинное).

И вот замок. \emph{Терминирует} (тогда это Element) или \emph{не терминирует} (тогда role-limit) ---
это ровно категорное \emph{вырезать} или \emph{слить}: статичная двойственность (ко)эквалайзеров
(§3.6) и динамическая (терминация или нет) --- одно и то же. Опоры рядом: Гильбертово пространство ---
процесс (завершённое пространство как объект $=\,\neg$P4); конституция --- условие на процесс
(статичное правило $=$ role-limit).\par}

\rezultat
Статика и динамика теории --- \textbf{две грани одной субстанции} с одной границей. «Вещественное»,
«бесконечномерное» суть \emph{процессы}, а не завершённые объекты; завершить их как объект ---
значит перейти границу.

% =============================================================================
\section*{7.3\quad Аудит стен: единственная несущая}
% =============================================================================

Расслоим все «стены» теории на три яруса. \emph{Ярус~I --- мнимые}: это \emph{черты}, а не дефекты
(нет нулевого объекта; 2-клетки Prop-тонкие). \emph{Ярус~II --- мягкие}: они \emph{растворены} ---
произведение для любой согласованной конституции, а не только эквивалентности; разрешимый фактор;
настоящая под-система. \emph{Ярус~III --- несущая}: ровно одна --- H1/role-limit $+$ L3 $=$ P4, и она
\emph{нерастворима by design}.

\rezultat
У теории ровно \textbf{одна} несущая граница. Всё прочее --- либо черта, либо растворённая мягкая
стена. Вот точный смысл «единственной стены»: не бедность теории, а её устройство.

% =============================================================================
\section*{7.4\quad Самоподобие}
% =============================================================================

Главный вывод всей арки. Граница H1 прежде жила в \emph{приложениях} --- в математике, в физике.
Теперь она установлена \emph{внутри собственного ядра}: двойственность (ко)эквалайзеров (§3.6),
незавершённая башня уровней (§6.2), аттрактор-как-role-limit (§4.5), топос-фрагмент (§7.1).

Теория \emph{демонстрирует} ту самую границу, которую предсказывает повсюду. Она \textbf{самоподобна}:
предел, который она находит в мире, она находит и в себе.

\rezultat
Единственная стена --- не внешнее ограничение, наложенное на теорию, а её \emph{внутренняя черта}.
Теория \textbf{самоподобна} границе H1/P4 --- и это, а не отдельная теорема, есть содержательный итог
блока.

% =============================================================================
\section*{7.5\quad Гранд-тур: вся теория в миниатюре}
% =============================================================================

Чтобы увидеть всё разом, проведём \emph{одну} простейшую систему --- дискретную, двухэлементную ---
через \emph{весь} аппарат: объект, оператор, 2-клетка, эквалайзер и коэквалайзер, образ, башня. На
одной этой системе видны обе ноты сразу --- позитивное начало (единица) и role-limit-граница.

\rezultat
Вся теория сворачивается в \textbf{миниатюру} на одной системе: H1 и P4 видны разом. Книга,
открывшаяся единицей (§1.2), здесь \emph{смыкает} обе ноты --- старт и предел --- в одной точке.

\medskip
\noindent\emph{Переход.} Теория замкнута на одной границе. Последний блок показывает, как из неё
\emph{ветвятся} миры --- математика, физика, знание, рассуждение.

% =============================================================================
\appendix
\section*{Приложение к Блоку 7 — машинные якоря}
% =============================================================================
\small\raggedright\sloppy
\noindent Всё перечисленное машинно проверено; единственная аксиома ветви --- \texttt{classic} (L3).
Конструкции стандартны; \emph{новое} --- развитие изнутри и самоподобие H1 (не новая теория категорий,
не «теория систем $=$ топос»).

\begin{itemize}\itemsep2pt
\item \textbf{§7.1.} классификатор подобъектов --- \texttt{char}, \texttt{char\_classifies},
\texttt{classified\_implies\_decidable} (H1), \texttt{classified\_implies\_saturated} (ярус Roles).\\
\emph{Файл:} \texttt{ERRSubobjectClassifier.v}.

\item \textbf{§7.2.} мост --- \texttt{dyn\_process}, \texttt{halve\_real} ($=$ вещественное-процесс),
\texttt{dyn\_collapse\_terminates} / \texttt{tower\_process\_not\_terminates}
(\texttt{ERRProcessBridge.v}); Гильберт-процесс \texttt{ERRHilbertProcess.v}; конституция-процесс
\texttt{ERRProcessConstitution.v}; \texttt{ERRProcess.v}.

\item \textbf{§7.3.} аудит стен --- свод §12; растворённые \texttt{ERRProductGeneral.v},
\texttt{ERRFiniteQuotient.v}, \texttt{ERRTierIIResidue.v}; H1-стратификация
\texttt{RoleLimitSpecies.v}.

\item \textbf{§7.4.} самоподобие --- синтез §3.6\,/\,§4.5\,/\,§6.2\,/\,§7.1 (H1 внутри ядра).

\item \textbf{§7.5.} гранд-тур одной системы через весь аппарат --- \texttt{ERRGrandTour.v}.
\end{itemize}
\normalsize

% --- Дальше: Блок 8 (ВЕТВЛЕНИЕ, закрытие) по тезисному плану в Книги/Теория Систем/. ---

\end{document}
